Volem construir un sensor de radiació ultraviolada que sigui sensible a radiacions de longitud d'ona de 300~nm. Decidim utilitzar l'efecte fotoelèctric com a principi del sensor. Així doncs, utilitzarem una cèl·lula fotoelèctrica que emeti electrons. Per al bon funcionament d'aquesta cèl·lula, cal que l'energia mínima dels electrons emesos sigui d'1~eV.

\begin{apartat}{1,25}
Calculeu la longitud d'ona llindar del material que hauríem d'utilitzar per a construir la cèl·lula.
\end{apartat}

\begin{apartat}{1,25}
Empleneu la taula de sota amb els valors de la longitud d'ona llindar dels tres materials donant el resultat en nanòmetres. Si podem triar un dels tres materials mostrats a la taula de sota per a construir la cèl·lula, quin triaríeu? Justifiqueu la resposta.
\end{apartat}

\begin{center}
\begin{tabular}{|c|c|c|c|}
\hline
\textit{Element} & \textit{Símbol} & \textit{Funció de treball (J)} & \textit{Longitud d'ona llindar (nm)} \\
\hline
tungstè & W & $8,36\times 10^{-19}$ & \\
\hline
magnesi & Mg & $5,86\times 10^{-19}$ & \\
\hline
potassi & K & $3,67\times 10^{-19}$ & \\
\hline
\end{tabular}
\end{center}

\dades{%
  $1\ \mathrm{eV} = 1,602\times 10^{-19}\ \mathrm{J}$.\\
  $c = 3,00\times 10^{8}\ \mathrm{m\,s^{-1}}$.\\
  $h = 6,63\times 10^{-34}\ \mathrm{J\,s}$.}
